\section{A threefold with an integral restriction kernel}\label{sec:injectivity}

We now construct a second threefold. Its role differs from that of the product: integral restriction itself fails in the first two degrees. All varieties in this section are over an algebraically closed field of characteristic two. A \emph{Tango datum} on a smooth projective variety $B$ means an effective ample divisor $D$ and a rational function $f$ such that $df$, divided locally by an equation of $2D$, is a regular nowhere-zero one-form. The stronger nowhere-zero condition permits iteration while preserving smoothness. We use the iteration of \citet[Proposition 2.3, p.~8 of the preprint]{Mukai}, but give the local construction in the quotient convention used throughout this paper.

\begin{lemma}[One iteration]\label{lem:iteration}
Suppose $(B,D,f)$ is a Tango datum and $D=qN$, where $q>2$ is odd and $N$ is effective. There are a smooth projective variety $\widetilde B$, a flat morphism $g:\widetilde B\to B$ with geometrically integral curve fibres, and a section $E$, with a Tango datum
\[
 \widetilde D=(q-1)E+g^*N.
\]
Let $0\to\OO_B\to V\to\OO_B(D)\to0$ be the nonzero Frobenius-killed extension defined by $df$, and set $P=\PP_B(V)$. The map $g$ factors through a finite flat map $\psi:\widetilde B\to P$ of degree $q$. For every integer $a$,
\begin{equation}\label{eq:iteration}
 \psi_*\OO_{\widetilde B}(aE)=
 \bigoplus_{i=0}^{q-1}\OO_P\!\left(\left\lfloor\frac{a-2i}{q}\right\rfloor\right)
             \otimes\pi^*\OO_B(2iN).
\end{equation}
For each prime $\ell\ne2$, pullback induces
$H^1_{\et}(B,\mathbf Q_\ell)\simeq H^1_{\et}(\widetilde B,\mathbf Q_\ell)$.
\end{lemma}
\begin{proof}
The Cartier exact-differential sequence, twisted by $\OO_B(-D)$, sends the normalized differential to a class in $H^1(B,\OO_B(-D))$ killed by Frobenius. It is nonzero because the preceding group $H^0(B,\OO_B(-2D))$ vanishes. This is the same construction as for the curve in Section~\ref{sec:surface}; in higher dimension the exact sequence is given in \citet[Theorem 1 and Lemma 1.4, pp.~1, 4--5 of the preprint]{Mukai}.

Let $b_i$ be local equations for $N$, and put $a_i=b_i^q$. Cartier gives, after restricting the open sets, regular functions $f_i$ and rational functions $h_i$ such that
\[
 f/a_i^2=f_i+h_i^2,\qquad df_i=a_i^{-2}df.
\]
The $df_i$ are nowhere zero. On overlaps write $t_{ij}=a_i/a_j$ and $c_{ij}=t_{ij}h_i-h_j$. Then
$f_j=t_{ij}^2f_i+c_{ij}^2$. The functions $c_{ij}$ are regular, since their squares are regular and the rings are normal. Glue the affine charts of $P$ by
$S_j=t_{ij}S_i+c_{ij}$. Their infinity section $E_0$ has normal bundle $\OO_B(D)$ and divisor bundle $\OO_P(1)$. The equations $S_i^2=f_i$ define a smooth divisor $G$ disjoint from $E_0$, with
\[
 \OO_P(G)=\OO_P(2)\otimes\pi^*\OO_B(-2D).
\]
Smoothness follows from the nowhere-zero $df_i$.

Set $M=\OO_P(-1)\otimes\pi^*\OO_B(2N)$. The relation
$M^{-q}=\OO_P(G+(q-2)E_0)$ defines a cyclic cover; take its normalization. Near $G$ the branch multiplicity is one. Near $E_0$, after an étale extension extracting a root of a unit, its equation is $w^q=s^{q-2}$, whose normalization has $s=u^q$, $w=u^{q-2}$. These are regular local models because $\gcd(q,q-2)=1$. They prove smoothness and $\psi^*E_0=qE$. The equation is Eisenstein at $G$, so the cover is integral. In eigendegree $i$, normalization and twisting by $aE$ allow the pole $\lfloor((q-2)i+a)/q\rfloor E_0$. Thus its eigensheaf is
\[
 M^i\bigl(\lfloor((q-2)i+a)/q\rfloor E_0\bigr),
\]
which is precisely \eqref{eq:iteration}. The algebra is locally free, so $\psi$ and $g$ are flat.

On the affine chart the normalized cover has equation
$T_i^q=S_i^2-f_i$. Its transition is $T_j=(b_i/b_j)^2T_i$, and hence $\widetilde f=b_i^2T_i$ is a global rational function. Since
\[
 d\widetilde f=b_i^2dT_i,\qquad qT_i^{q-1}dT_i=-df_i,
\]
the normalized differential is nonzero off the branch. At $G$, one can use $f_i$ as a base coordinate and solve the equation for it; $T_i$ is then a smooth parameter and $dT_i$ is nonzero. At $E$, the function $T_i$ has pole order two, while $df_i$ pulls back to a nowhere-zero base differential. The displayed identity therefore gives exact zero order $2(q-1)$ for $dT_i$. This proves the asserted Tango condition for $\widetilde D$.

A geometric fibre has affine equation $T^q=(S-c)^2$ and one point at infinity. Its normalization is $\PP^1$, and that normalization is a universal homeomorphism. The fibre is integral, and $E$ has degree one on it. Moreover
$(\OO(E)|_E)^q\simeq\OO_B(qN)$, so $E|_E$ is numerically ample. The effective divisor $E$ is nef: curves outside $E$ intersect it nonnegatively, and its restriction to $E$ is ample. For a $d$-dimensional subvariety $Z$ of $\widetilde B$, the image in $B$ has dimension $d$ or $d-1$. In the first case $(g^*N)^d\cdot Z>0$; in the second, $E(g^*N)^{d-1}\cdot Z>0$. All other terms in $\widetilde D^d\cdot Z$ are nonnegative. Nakai--Moishezon proves ampleness.

Finally each geometric fibre has the étale cohomology of $\PP^1$. Proper base change gives $g_*\mathbf Q_\ell=\mathbf Q_\ell$ and $R^1g_*\mathbf Q_\ell=0$. The degree-one Leray sequence proves the last assertion. These uses of proper base change may be applied with $\mathbf Z/\ell^n$ first and then passed to $\mathbf Q_\ell$ \citep[Exposé XII, Theorem 5.1 and Corollary 5.2, pp.~88--89]{SGA4}.
\end{proof}

\begin{definition}[The second threefold]\label{def:threefold}
Start from the smooth completion of $C:y^2+y=x^{843}$. Let $P_\infty$ be its point at infinity and $N_0=\OO_C(2P_\infty)$. The function $x$ has
$(dx)=840P_\infty$, so $D_1=420P_\infty=210N_0$ is a Tango divisor. Apply Lemma~\ref{lem:iteration} with $q=21$ and $N=10N_0$, obtaining $f:S\to C$ and $E_1$. Put
\[
 B=\OO_S(2E_1)\otimes f^*N_0.
\]
The new Tango divisor is $10B$. Apply the lemma again with $q=5$ and $N=2B$, obtaining $g:R\to S$ and $E_2$, and put
\[
 L=\OO_R(2E_2)\otimes g^*B.
\]
Its new Tango divisor is $2L$. Choose 843 general sections of $L^{445}$, with the first eight defining a closed immersion in $\PP^7$, and take their simultaneous 445th-root cover $\rho:U\to R$. Put $H=\rho^*L$.
\end{definition}

We prove all the existence and positivity assertions implicit in this definition before using it. Letters $U,H$ distinguish this threefold from the product $Z,B$ of Theorem~\ref{thm:main}; within this section $B$ denotes the surface polarization in Definition~\ref{def:threefold}.

\begin{proposition}[Positivity and first Betti number]\label{prop:threefold-positive}
The curve in Definition~\ref{def:threefold} has genus $421$. The line bundles $B$ and $L$ are ample and generated, $B^{441}$ and $L^{445}$ are very ample, and
\[
 h^0(S,B)=24,\qquad h^0(S,B^5)=336.
\]
The root cover $U$ is smooth, projective, and connected, $H$ is very ample, and
\[
 \deg\rho=445^{843},\qquad
 b_1(U):=\dim H^1_{\et}(U,\mathbf Q_3)=842.
\]
\end{proposition}
\begin{proof}
The calculation in Lemma~\ref{lem:curve}, with $h=422$, gives $g(C)=421$ and the stated divisor of $dx$. The sections $1,x$ generate $N_0$. Formula~\eqref{eq:iteration} gives
\[
 f_*B=N_0\oplus N_0^{21},\qquad g_*L=B\oplus B^5.
\]
In both steps the index-zero sections generate away from the ramification section; the index-one local generator $w/s=u^{-2}$ is nonvanishing in the frame of $\OO(2E)$ on that section. The coefficient bundles are generated. Hence $B$ and $L$ are generated. They are ample because $10B$ and $2L$ are the ample divisors of the iteration.

The extension on the first projective bundle is
$0\to\OO_C\to V_1\to N_0^{210}\to0$. For $a\ge1$ and $b\ge441$, the bundle $\Sym^aV_1\otimes N_0^b$ has a filtration with quotients $N_0^{b+210j}$. Each has degree at least $882>2g(C)+1$. Serre duality shows that its $H^1$ vanishes after subtracting every length-two subscheme. Sections therefore provide all base values and first jets, while $\OO(a)$ provides the complete fibre system. Thus $\OO_{\PP(V_1)}(a)\otimes N_0^b$ is very ample.

Now $B^{441}=\psi_1^*F_1$, where $F_1=\OO(42)\otimes N_0^{441}$. The cover algebra has summands
\[
 K_i=\OO\!\left(-\left\lceil2i/21\right\rceil\right)\otimes N_0^{20i},
 \quad 0\le i\le20.
\]
Every $F_1K_i$ is generated: its relative degree is at least 40 and its base exponent at least 441. Sections from $F_1$ recover the embedded projective bundle, and generators from all $F_1K_i$ recover its finite cover algebra on each affine chart. This proves very ampleness of $B^{441}$. More generally, for $t\ge441$ the eigensheaves of $B^t$ push to
\[
 \Sym^{\lfloor2(t-i)/21\rfloor}V_1\otimes N_0^{t+20i}.
\]
They have nonnegative relative degree and filtrations with line quotients of degree at least $882>2g-2$. Consequently $H^1(S,B^t)=0$. Since $B^t$ is very ample, its evaluation on every length-two subscheme $Z$ is surjective, and
\begin{equation}\label{eq:B-jets}
 H^1(S,B^t\otimes I_Z)=0\quad(t\ge441,\ \operatorname{length}Z\le2).
\end{equation}

The second extension is $0\to\OO_S\to V_2\to B^{10}\to0$. By \eqref{eq:B-jets}, the filtered bundle $\Sym^aV_2\otimes B^b$ evaluates surjectively on every length-two subscheme of $S$ whenever $a\ge1$, $b\ge441$. Fibre sections separate the remaining vertical directions. Hence $\OO_{\PP(V_2)}(a)\otimes B^b$ is very ample in this range. Since
\[
 L^{445}=\psi_2^*\bigl(\OO(178)\otimes B^{445}\bigr),
\]
and the normalized algebra summands are
$\OO(-\lceil2j/5\rceil)\otimes B^{4j}$, all their twists by $\OO(178)\otimes B^{445}$ are generated. The same finite-algebra argument proves that $L^{445}$ is very ample.

For $0\le b\le105$, the pole semigroup of $C$ gives $h^0(N_0^b)=b+1$: only $1,x,\ldots,x^b$ have pole order at most $2b<843$. Thus
\[
 h^0(B)=2+22=24,\qquad
 h^0(B^5)=\sum_{i=0}^5(6+20i)=336.
\]
There are at least 179 independent sections of $L^{445}$, since its index-zero pushforward has 179 positive line quotients with vanishing $H^1$. A general projection of the resulting threefold embedding to $\PP^7$ remains a closed immersion: the closed secant and tangent locus has dimension at most seven, and a general centre of codimension eight misses it. Thus eight sections can be chosen to embed $R$. Independently, 843 general divisors are smooth and meet with simple normal crossings, with no fourfold intersection. Both requirements are nonempty open conditions on an irreducible parameter space, so they hold together.

The root-cover local coordinates prove smoothness exactly as in Section~\ref{sec:roots}, now with at most three branch parameters. Its algebra is locally free of rank $445^{843}$. The decomposition into negative powers of the ample $L$ gives $H^0(U,\OO_U)=k$, so it is connected. A smooth connected variety is integral. The generated sections of $L$ together with the root sections embed $U$: on a chart $s\ne0$ of $L$, the functions $b_i/s^{445}$ from the first eight sections generate the base affine algebra after adjoining the section $s^{445}$ to their very ample system, and $(z_i/s)^{445}=b_i/s^{445}$. The ratios $z_i/s$ consequently generate the finite cover algebra. Properness finishes the closed-immersion assertion.

The two Tango steps preserve $H^1_{\et}$, so $b_1(R)=842$. To see that the root cover preserves this number, extend $\mathbf Q_3$ to a finite extension $K$ containing all 445th roots of unity and decompose $\rho_*K$ under $(\mu_{445})^{843}$. Its trivial character is $K_R$. For a nontrivial character $\chi$, let $J$ be the nonempty set of its nontrivial inertia indices. The corresponding sheaf is $j_!\mathcal L_\chi$, where $j:V=R\setminus\bigcup_{i\in J}D_i\hookrightarrow R$ and $\mathcal L_\chi$ is lisse of rank one. Locally this is the character decomposition of each tame root cover: a nontrivial character has zero stalk on its branch divisor, and a trivial one extends constantly. The open $V$ is smooth affine of dimension three because the removed divisor is ample. Affine vanishing and Poincaré duality give
\[
 H^1(R,j_!\mathcal L_\chi)=H_c^1(V,\mathcal L_\chi)
   =H^5(V,\mathcal L_\chi^\vee(3))^\vee=0.
\]
These are \citet[Exposé XIV, Theorem 3.1 and Corollary 3.2, p.~160; Exposé XVIII, Theorem 3.2.5 and (3.2.6.2), pp.~585--587]{SGA4}. Finite proper base change has no higher direct images for $\rho$. Thus only the trivial character contributes in degree one, proving $b_1(U)=842$.
\end{proof}

\subsection{The exact negative cohomology}

\begin{lemma}\label{lem:threefold-negative}
For the line bundles in Definition~\ref{def:threefold},
\[
 h^1(R,L^{-r})=\begin{cases}360&r=1,\\1&r=2,\\0&r\ge3,\end{cases}
 \qquad
 h^1(U,H^{-d})=\begin{cases}1203&d=1,\\1&d=2,\\0&d\ge3.\end{cases}
\]
\end{lemma}
\begin{proof}
Let $\alpha\in H^1(S,B^{-10})$ be the extension class of $V=V_2$. For $0\le e<10$,
\begin{equation}\label{eq:nonannihilation}
 H^0(S,V^\vee B^e)=0.
\end{equation}
Indeed this group is the kernel of multiplication by $\alpha$ from $H^0(B^e)$ to $H^1(B^{e-10})$. Under the Frobenius-kernel description, multiplying $\alpha$ by a nonzero section $s$ replaces its rational exact differential $\eta$ by $s^2\eta$. This remains nonzero. The connecting map for the ample divisor $(10-e)B$ is injective, since $H^0(B^{-2(10-e)})=0$. Thus $s\alpha\ne0$, proving \eqref{eq:nonannihilation}.

Put $m_j(r)=\lceil2(r+j)/5\rceil$ and $e_j(r)=4j-r-10$. Formula~\eqref{eq:iteration} and relative duality give
\begin{equation}\label{eq:threefold-negative}
 H^1(R,L^{-r})=
 \bigoplus_{m_j(r)\ge2}
 H^0\bigl(S,(\Sym^{m_j(r)-2}V)^\vee B^{e_j(r)}\bigr).
\end{equation}
All zeroth direct images vanish. The dual in this formula is taken after the symmetric power; commuting these operations is not valid in characteristic two.

For $r\ge7$, every $e_j(r)\le6-r<0$, and the dual symmetric power has only quotients $B^{-10a}$ in its filtration, so all sections vanish. At $r=3$, the only possible term is $H^0(V^\vee B^3)$, zero by \eqref{eq:nonannihilation}. At $r=4,5,6$, the only possible term is $H^0((\Sym^2V)^\vee B^{6-r})$. Frobenius splitting gives the decomposition
\begin{equation}\label{eq:sym-two}
 \Sym^2V\simeq V\oplus B^{20}.
\end{equation}
For clarity, if $u$ generates the trivial subline and $v_j=a_{ij}v_i+c_{ij}u$, the subbundle $uV$ is generated by $u^2,uv_i$ and is isomorphic to $V$. A splitting of $F^*V$ provides combinations $v_i^2+h_i u^2$ that glue as $B^{20}$. They complement $uV$. This proves \eqref{eq:sym-two}; its dual and \eqref{eq:nonannihilation} give the required vanishing.

At $r=2$, the $j=3$ term in \eqref{eq:threefold-negative} is $H^0(\OO_S)=k$, while the $j=4$ term is $H^0(V^\vee B^4)=0$. At $r=1$, the only nonzero terms are $H^0(B)$ and $H^0(B^5)$, of dimensions 24 and 336. This proves the first formula. Finally
\[
 H^1(U,H^{-d})=
 \bigoplus_{0\le a_i<445}H^1(R,L^{-(d+\sum a_i)}).
\]
There are 843 tuples with sum one. Hence the dimensions are $360+843=1203$, $1$, and $0$ in the asserted ranges.
\end{proof}

\subsection{Negative one-forms}

The second iteration supplies a nowhere-zero map $\vartheta:L^4\to\Omega_R^1$. Use the same letter for its pullback $H^4\to\Omega_U^1$; the latter is generically nonzero because the cover is separable.

\begin{lemma}\label{lem:threefold-forms}
For every $r\ge2$ and $d\ge2$,
\[
 H^0(R,\Omega_R^1L^{-r})=\vartheta H^0(R,L^{4-r}),\qquad
 H^0(U,\Omega_U^1H^{-d})=\vartheta H^0(U,H^{4-d}).
\]
The map $\vartheta:H^4\to\Omega_U^1$ does not vanish identically on any smooth divisor $Y\in|H^d|$, $d\ge1$.
\end{lemma}
\begin{proof}
The first iteration gives a subbundle $B^{20}\subset\Omega_S^1$ with line-bundle quotient $Q$. On $P=\PP_S(V)$, the logarithmic cotangent bundle for the disjoint branches $E_0,G$ has a filtration with quotients
\[
 \OO_P(2),\qquad \pi^*Q,\qquad \OO_P(-1)\otimes B^{10}.
\]
To check this locally, write $G$ as $S^2-f_i=0$, with $df_i$ spanning the distinguished base line. The horizontal logarithmic generator $df_i/(S^2-f_i)$ allows precisely the twist $B^{20}(G)=\OO_P(2)$; the remaining horizontal quotient is $Q$. The relative logarithmic quotient has a pole only at $E_0$, hence is $\Omega^1_{P/S}(E_0)=\OO_P(-1)B^{10}$. In the chart $z=1/S$, these generators become $df_i/(1-f_i z^2)$ and $dz/z$, giving the same filtration across infinity.

The tame normalized cover embeds $\Omega_R^1$ in the pullback of this logarithmic bundle. After twisting by $L^{-r}$, the last two pieces have relative degrees $-m_j(r)$ and $-1-m_j(r)$, and therefore no global sections. The first piece has eigensheaves
\[
 \OO_P(2-m_j(r))\otimes B^{4j-r}.
\]
For $r\ge6$ all relative degrees are negative. At $r=5$, the only degree-zero term is $B^{-5}$, which also has no sections. At $r=4$, the only nonzero possibility is $H^0(\OO_S)$. At $r=3$ the possibilities are $H^0(B)\oplus H^0(B^5)=H^0(R,L)$. At $r=2$ they are
\[
 H^0(VB^{-2})\oplus H^0(B^2)\oplus H^0(B^6)\oplus H^0(B^{10}).
\]
Here $H^0(VB^{-2})=0$: its extension sequence has connecting map $H^0(B^8)\to H^1(B^{-2})$, and \eqref{eq:nonannihilation} proves its injectivity. The other three terms equal $H^0(R,L^2)$ by \eqref{eq:iteration}. Thus in every case the logarithmic calculation bounds the left side by $h^0(L^{4-r})$. Multiplication by the nonzero $\vartheta$ injects that latter space into it, proving equality.

Let $D$ be the branch divisor of $\rho$. The residue sequence gives
\[
 H^0(R,\Omega_R^1(\log D)L^{-r})=H^0(R,\Omega_R^1L^{-r})
 \quad(r>0),
\]
since $L^{-r}|_{D_i}$ is antiample. The root-cover logarithmic inclusion and its eigensheaves therefore bound
\[
 h^0(U,\Omega_U^1H^{-d})
 \le\sum_{0\le a_i<445}h^0(R,L^{4-d-\sum a_i})=h^0(U,H^{4-d}).
\]
The pulled-back $\vartheta$ provides the opposite injection, and equality follows. If $\vartheta$ vanished identically along a divisor in $|H^d|$, its defining section of $\Omega_U^1H^{-4}$ would come from $H^0(\Omega_U^1H^{-4-d})$, which vanishes by the formula. This proves the last assertion without claiming that the pulled-back map is everywhere a subbundle.
\end{proof}

\subsection{The integral lattice in degree one}

\begin{lemma}\label{lem:first-lattice}
Let $h:U\to C$ be the composite morphism. Then $h_*\OO_U=\OO_C$, its geometric generic fibre is integral, and pullback gives
\[
 C^1(C)\xrightarrow{\sim}C^1(U).
\]
For every smooth $Y\in|H^d|$ with $d\ge2$, one also has $(h|_Y)_*\OO_Y=\OO_C$ and integral isomorphisms
\[
 C^1(C)\xrightarrow{\sim}C^1(Y),\qquad C^1(U)\xrightarrow{\sim}C^1(Y).
\]
\end{lemma}
\begin{proof}
Both Tango maps are flat with geometrically integral fibres. Their composite $R\to C$ has the same property: on a geometric fibre the second map is flat over the integral first fibre, so affine charts inject into their integral generic-fibre algebras, proving reducedness and irreducibility. Thus $k(R)/k(C)$ is regular. Let $U\to C'\to C$ be the Stein factorization. The curve $C'$ is smooth, and regularity implies that $k(R)$ and $k(C')$ are linearly disjoint over $k(C)$. Therefore
\[
 [k(C'):k(C)]\mid[k(U):k(R)]=445^{843}.
\]
In particular $C'\to C$ is separable. The map $U\to C'$ has geometrically connected fibres, so $H^1_{\et}(C',\mathbf Q_3)$ injects into $H^1_{\et}(U,\mathbf Q_3)$. Hence $g(C')\le421$. Riemann--Hurwitz gives
$g(C')-1\ge\deg(C'/C)(421-1)$, forcing degree one. This proves $h_*\OO_U=\OO_C$. The extension $k(U)/k(C)$ is separably generated, since $k(R)/k(C)$ is regular and the root cover is separable. It has no new algebraic constants, so it is regular and the geometric generic fibre is integral.

The two Tango maps have sections, and their pullbacks on reduced Picard varieties are split monomorphisms. Their first Betti numbers coincide, so these maps are isomorphisms. Thus $\operatorname{Pic}^0_{R,\mathrm{red}}\simeq\operatorname{Pic}^0_C$. For the finite locally free map $\rho$, the norm identity
\[
 \operatorname{Nm}_\rho\circ\rho^*=[445^{843}]
\]
shows that pullback on reduced Picard varieties is an isogeny with kernel killed by an odd integer. Here the dimensions agree by Proposition~\ref{prop:threefold-positive}. Its map on 2-divisible groups is therefore an isomorphism. The integral Dieudonné realization of the reduced Picard variety identifies these groups functorially with first crystalline cohomology \citep[II, Proposition 3.11 and Remark 3.11.2, pp.~618--619]{Illusie}. This proves the integral assertion for $U$.

Now let $d\ge2$. The divisor $Y$ dominates $C$: a divisor supported on finitely many fibres would miss another positive-dimensional fibre and could not be ample. The geometric generic fibre $U_{\overline\eta}$ is an integral projective Cohen--Macaulay surface, since it is a geometric fibre of a flat map from a smooth variety to a smooth curve. The restriction $Y_{\overline\eta}$ is an ample effective Cartier divisor. It is connected. Indeed, Serre duality and Serre vanishing give $H^1(\OO_{U_{\overline\eta}}(-nY))=0$ for large $n$, and the divisor sequence then gives $H^0(\OO_{nY_{\overline\eta}})=\overline{k(C)}$. A disconnected $Y_{\overline\eta}$ would yield a nontrivial idempotent on every $nY_{\overline\eta}$, contradicting this equality.

It follows that the finite Stein curve map $C_Y\to C$ is purely inseparable. We rule out a nontrivial such map by keeping track of the coefficient of the Tango differential. Write $h_R:R\to C$, and let $G_1'\subset S$ and $G_2'\subset R$ be the reduced inverse images of the two inseparable multisections. The effective divisor
\[
 D_0=210h_R^*P_\infty+5g^*G_1'+G_2'
\]
has line bundle $L$. Indeed,
\[
 \OO_S(G_1')=\OO_S(2E_1)\otimes f^*N_0^{-20},\qquad
 \OO_R(G_2')=\OO_R(2E_2)\otimes g^*B^{-4},
\]
so the first two summands of $D_0$ have line bundle $g^*B^5$, and the last changes it to $L$. Let $\sigma_0\in H^0(R,L)$ have divisor $D_0$. The local equations of the two covers give vanishing orders twenty and four for the successive pullbacks of the normalized base differential. Together with $(dx)=840P_\infty$, this gives
\[
 h_R^*dx=c\,\vartheta_R(\sigma_0^4),\qquad c\in k^\times,
\]
where $\vartheta_R:L^4\to\Omega_R^1$ is the second Tango subbundle. Both coefficients have divisor $4D_0$, so their ratio is a nonzero constant. Pullback gives the same identity on $U$.

The divisor $\rho^*D_0\in|H|$ cannot contain $Y\in|H^d|$ when $d\ge2$, because otherwise $H^{1-d}$ would have a nonzero section. Thus the coefficient of $\vartheta$ is nonzero at the generic point of $Y$. Moreover the image of $\vartheta$ in $\Omega_Y^1$ is nonzero: if it vanished, its restriction would factor through the conormal bundle $H^{-d}|_Y$, whereas $H^0(Y,H^{-d-4}|_Y)=0$; the absence of divisorial zeros of $\vartheta$ excludes a zero restriction. Therefore $(h|_Y)^*dx\ne0$, so $Y\to C$ is generically separable. The Stein map is both separable and purely inseparable and hence has degree one. This proves $C_Y=C$.

The equality $(h|_Y)_*\OO_Y=\OO_C$ makes $\operatorname{Pic}_C\to\operatorname{Pic}_Y$ a monomorphism of fppf Picard sheaves. To see this, the equality of structure-sheaf pushforwards persists after every base change over $k$, by flat base change. If a pulled-back line bundle is trivial, or comes from the parameter base, projection formula and pushforward recover the same assertion on $C$. After local fppf trivialization this proves the assertion for the sheaf quotient defining the Picard functor. Thus $\operatorname{Pic}^0_C\to\operatorname{Pic}^0_{Y,\mathrm{red}}$ is a closed immersion. Étale weak Lefschetz gives $b_1(Y)=b_1(U)=842$; equivalently, affine vanishing and duality on the smooth affine complement give the required degree-one isomorphism. The two abelian varieties both have dimension 421, so their closed immersion is an isomorphism. Their integral Dieudonné realizations give $C^1(C)\simeq C^1(Y)$. Factoring this map through $U$ proves the final assertion.
\end{proof}

\subsection{The exact injectivity range}

\begin{theorem}\label{thm:injectivity-exact}
For $(U,H)$ in Definition~\ref{def:threefold} and every smooth $Y_d\in|H^d|$, restriction
\[
 H^2_{\cris}(U/W)\longrightarrow H^2_{\cris}(Y_d/W)
\]
is injective if and only if $d\ge3$. In degree one its kernel killed by two has dimension at least 361, and in degree two this dimension is at least one. The coherent restriction kernels in $H^1(\OO)$ have dimensions $1203$, $1$, and $0$ for $d=1$, $d=2$, and $d\ge3$, respectively.
\end{theorem}
\begin{proof}
The coherent dimensions follow from Lemma~\ref{lem:threefold-negative}, the divisor sequence, and connectedness of each smooth ample section. The functorial conjugate edge injects this coherent kernel, after Frobenius twist, into the first de Rham restriction kernel. When $d=1$, its dimension is 1203. In the universal-coefficient sequence
\[
 0\to C^1(U)/2\to H^1_{\dR}(U/k)\to C^2(U)[2]\to0,
\]
the first term has dimension $842$. Naturality therefore supplies a subspace of dimension at least $1203-842=361$ in the kernel on $C^2[2]$.

For $d=2$, the coherent kernel is nonzero. In the natural diagram of the same universal-coefficient sequences for $U$ and $Y_2$, the left vertical map is an isomorphism by Lemma~\ref{lem:first-lattice}. The snake lemma identifies the first de Rham restriction kernel with the kernel on $C^2[2]$. The former contains the nonzero coherent kernel, proving failure of integral injectivity.

Let $d\ge3$, and first consider a global closed one-form $\omega$ on $U$ whose restriction to $Y=Y_d$ is zero. The conormal sequence has $H^0(Y,H^{-d}|_Y)=0$. Thus $\omega$ lies in the image of $H^0(\Omega_U^1H^{-d})$, multiplied by the defining section $s_Y$. For $d\ge5$ this space is zero by Lemma~\ref{lem:threefold-forms}. For $d=3,4$ write
\[
 \omega=s_Yt\vartheta,\qquad t\in H^0(U,H^{4-d}).
\]
If $t\ne0$, it cannot vanish on $Y$, since that would give a nonzero section of the antiample bundle $H^{4-2d}$. Locally write $s_Y=g$, $t=u$, and let $\vartheta_0$ represent $\vartheta$. Reduction of $d\omega=0$ modulo $g$ gives
$u\,dg\wedge\vartheta_0=0$. At the generic point of the integral smooth divisor $Y$, this implies that $\vartheta|_Y$ maps to zero in $\Omega_Y^1$. Since this latter bundle is torsion-free, the vanishing holds everywhere, and $\vartheta|_Y$ factors through the conormal line:
\[
 H^4|_Y\longrightarrow H^{-d}|_Y.
\]
There is no such nonzero homomorphism because $H^0(Y,H^{-d-4}|_Y)=0$. Thus $\vartheta$ itself would vanish along $Y$, contrary to Lemma~\ref{lem:threefold-forms}. Hence $\omega=0$.

For a smooth proper connected variety the Hodge edge identifies the kernel of $H^1_{\dR}\to H^1(\OO)$ with global closed one-forms; those forms inject into de Rham cohomology because the differential of every global function is zero. A de Rham restriction-kernel class has zero coherent edge, as coherent restriction is injective for $d\ge3$. It therefore comes from a closed global one-form restricting to zero, which the preceding paragraph proves is zero. Thus first de Rham restriction is injective.

The first crystalline lattices are isomorphic by Lemma~\ref{lem:first-lattice}. The universal-coefficient diagram now gives injectivity on $C^2[2]$. Rational crystalline weak Lefschetz makes the full kernel a finite torsion $W$-module. Every nonzero such module contains a nonzero element killed by two, so this kernel is zero. This proves the exact range.
\end{proof}
